% GENERATED FILE. Edit canonical lessons, then run npm run generate:books.
% canonical-source-hash: fnv1a64:1b163b61e324add6
% canonical-lessons: BN-C185-bori, BN-C185-injekshon, BN-C185-chot, BN-C185-phoska, BN-C185-bndej

\chapter{Pill, Injection, Injury, Blister, Bandage}
\label{ch:bn-pill-injection-injury-blister-bandage}
\hlchaptermodality{185}

\begin{chapteropening}
\textbf{By the end of this chapter:} \emph{I can say a pill, an injection, an injury, a blister and a bandage.}

Complete the last lesson of chapter 185: I can say a pill, an injection, an injury, a blister and a bandage.
\end{chapteropening}

\section[\texorpdfstring{\bn{বড়ি} (bôṛi)}{বড়ি (bôṛi)}]{\bn{বড়ি} (bôṛi) — a pill}

\label{lesson:BN-C185-bori}



\begin{quote}
Before the new one: say the Bengali for to obstruct, then the Bengali for to take a risk.
\end{quote}

\subsection*{You'll want to know: \bn{বড়ি}}
\textbf{\bn{বড়ি}} — \emph{bôṛi} — \textquotedblleft{}a pill\textquotedblright{}.

\begin{cousinweb}
\textbf{\bn{বড়ি}} (\emph{bôṛi}) is a tadbhava, a word inherited from Sanskrit and reshaped by sound change. It comes from \emph{vaṭa}, \textquotedblleft{}small lump; round cake of pulse\textquotedblright{} (CDIAL 11213). CDIAL lists Bengali \emph{baṛi}, \textquotedblleft{}ball of bruised pulse, pill\textquotedblright{}.
\end{cousinweb}

\begin{grammarlens}[title={What you've built}]
One more word for this chapter.
\end{grammarlens}

\subsection*{Your turn}
\begin{itemize}
\raggedright
  \item \emph{Say it:} \emph{bôṛi}
  \item \emph{Say it:} \emph{bôṛi}, once more
  \item \emph{Say it:} say \emph{bôṛi}
  \item \emph{Recall:} say the Bengali for to obstruct, then the Bengali for to take a risk, then say \emph{bôṛi} again
\end{itemize}

\begin{tcolorbox}[breakable,colback=teal!4,colframe=teal!35!black,title={Before you move on}]
What does \bn{বড়ি} mean? (\textquotedblleft{}A pill\textquotedblright{}.) Say it once more.
\end{tcolorbox}
\section[\texorpdfstring{\bn{ইনজেকশন} (injekshôn)}{ইনজেকশন (injekshôn)}]{\bn{ইনজেকশন} (injekshôn) — an injection}

\label{lesson:BN-C185-injekshon}



\begin{quote}
Before the new one: say the Bengali for to take a risk, then the Bengali for a pill.
\end{quote}

\subsection*{You'll want to know: \bn{ইনজেকশন}}
\textbf{\bn{ইনজেকশন}} — \emph{injekshôn} — \textquotedblleft{}an injection\textquotedblright{}.

\begin{cousinweb}
\textbf{\bn{ইনজেকশন}} (\emph{injekshôn}) is a loan from English \emph{injection}.
\end{cousinweb}

\begin{grammarlens}[title={What you've built}]
One more word for this chapter.
\end{grammarlens}

\subsection*{Your turn}
\begin{itemize}
\raggedright
  \item \emph{Say it:} \emph{injekshôn}
  \item \emph{Say it:} \emph{injekshôn}, once more
  \item \emph{Say it:} say \emph{injekshôn}
  \item \emph{Recall:} say the Bengali for to take a risk, then the Bengali for a pill, then say \emph{injekshôn} again
\end{itemize}

\begin{tcolorbox}[breakable,colback=teal!4,colframe=teal!35!black,title={Before you move on}]
What does \bn{ইনজেকশন} mean? (\textquotedblleft{}An injection\textquotedblright{}.) Say it once more.
\end{tcolorbox}
\section[\texorpdfstring{\bn{চোট} (choṭ)}{চোট (choṭ)}]{\bn{চোট} (choṭ) — an injury}

\label{lesson:BN-C185-chot}



\begin{quote}
Before the new one: say the Bengali for a pill, then the Bengali for an injection.
\end{quote}

\subsection*{You'll want to know: \bn{চোট}}
\textbf{\bn{চোট}} — \emph{choṭ} — \textquotedblleft{}an injury\textquotedblright{}.

\begin{cousinweb}
\textbf{\bn{চোট}} (\emph{choṭ}) is a tadbhava, a word inherited from Sanskrit and reshaped by sound change. It comes from \emph{cuṭati}, \textquotedblleft{}cuts off, strikes\textquotedblright{} (CDIAL 4857). CDIAL lists Bengali \emph{coṭ}, \textquotedblleft{}to hit\textquotedblright{}, and calls the group non-Aryan.
\end{cousinweb}

\begin{grammarlens}[title={What you've built}]
One more word for this chapter.
\end{grammarlens}

\subsection*{Your turn}
\begin{itemize}
\raggedright
  \item \emph{Say it:} \emph{choṭ}
  \item \emph{Say it:} \emph{choṭ}, once more
  \item \emph{Say it:} say \emph{choṭ}
  \item \emph{Recall:} say the Bengali for a pill, then the Bengali for an injection, then say \emph{choṭ} again
\end{itemize}

\begin{tcolorbox}[breakable,colback=teal!4,colframe=teal!35!black,title={Before you move on}]
What does \bn{চোট} mean? (\textquotedblleft{}An injury\textquotedblright{}.) Say it once more.
\end{tcolorbox}
\section[\texorpdfstring{\bn{ফোসকা} (phoskā)}{ফোসকা (phoskā)}]{\bn{ফোসকা} (phoskā) — a blister}

\label{lesson:BN-C185-phoska}



\begin{quote}
Before the new one: say the Bengali for an injection, then the Bengali for an injury.
\end{quote}

\subsection*{You'll want to know: \bn{ফোসকা}}
\textbf{\bn{ফোসকা}} — \emph{phoskā} — \textquotedblleft{}a blister\textquotedblright{}.

\begin{cousinweb}
\textbf{\bn{ফোসকা}} (\emph{phoskā}) is a tadbhava, a word inherited and reshaped by sound change. It comes from \emph{*phossa}, \textquotedblleft{}hollow\textquotedblright{}, a form CDIAL reconstructs from its descendants (CDIAL 9114). CDIAL lists Bengali \emph{phoskā}, \textquotedblleft{}blister\textquotedblright{}.
\end{cousinweb}

\begin{grammarlens}[title={What you've built}]
One more word for this chapter.
\end{grammarlens}

\subsection*{Your turn}
\begin{itemize}
\raggedright
  \item \emph{Say it:} \emph{phoskā}
  \item \emph{Say it:} \emph{phoskā}, once more
  \item \emph{Say it:} say \emph{phoskā}
  \item \emph{Recall:} say the Bengali for an injection, then the Bengali for an injury, then say \emph{phoskā} again
\end{itemize}

\begin{tcolorbox}[breakable,colback=teal!4,colframe=teal!35!black,title={Before you move on}]
What does \bn{ফোসকা} mean? (\textquotedblleft{}A blister\textquotedblright{}.) Say it once more.
\end{tcolorbox}
\section[\texorpdfstring{\bn{ব্যান্ডেজ} (bænḍej)}{ব্যান্ডেজ (bænḍej)}]{\bn{ব্যান্ডেজ} (bænḍej) — a bandage}

\label{lesson:BN-C185-bndej}



\begin{quote}
Before the new one: say the Bengali for an injury, then the Bengali for a blister.
\end{quote}

\subsection*{You'll want to know: \bn{ব্যান্ডেজ}}
\textbf{\bn{ব্যান্ডেজ}} — \emph{bænḍej} — \textquotedblleft{}a bandage\textquotedblright{}.

\begin{cousinweb}
\textbf{\bn{ব্যান্ডেজ}} (\emph{bænḍej}) is a loan from English \emph{bandage}.
\end{cousinweb}

\begin{grammarlens}[title={What you've built}]
One more word for this chapter.
\end{grammarlens}

\subsection*{Your turn}
\begin{itemize}
\raggedright
  \item \emph{Say it:} \emph{bænḍej}
  \item \emph{Say it:} \emph{bænḍej}, once more
  \item \emph{Say it:} say \emph{bænḍej}
  \item \emph{Recall:} say the Bengali for an injury, then the Bengali for a blister, then say \emph{bænḍej} again
\end{itemize}

\begin{tcolorbox}[breakable,colback=teal!4,colframe=teal!35!black,title={Before you move on}]
What does \bn{ব্যান্ডেজ} mean? (\textquotedblleft{}A bandage\textquotedblright{}.) Say it once more.
\end{tcolorbox}
